% 分野別類題: 重積分 解答例
% コンパイル: TEXINPUTS="../../../00_common:$TEXINPUTS" latexmk -xelatex answers.tex
\documentclass[a4paper,11pt]{article}
\usepackage{preamble}
\usepackage{shoakubox}

\begin{document}

\kamokumei{類題演習 解答例 --- 重積分}

\daimon{【解法】重積分の解き方と導出}

\noindent
\textbf{(1) 累次積分と積分順序の交換（Fubini の定理）}

\noindent
連続関数 $f(x,y)$ が有界閉領域 $D$ 上で定義されているとき、$D$ を
\[
D = \{(x,y) \mid a \le x \le b,\ \varphi_1(x) \le y \le \varphi_2(x)\}
\]
のように $x$ を先に固定して $y$ の範囲を表す「縦線集合」として書ける場合、重積分は
\[
\iint_{D} f(x,y)\,dA = \int_{a}^{b}\left(\int_{\varphi_1(x)}^{\varphi_2(x)} f(x,y)\,dy\right)dx
\]
と累次積分に書き換えられる（Fubini の定理）。同じ領域 $D$ を「横線集合」
\[
D = \{(x,y) \mid c \le y \le d,\ \psi_1(y) \le x \le \psi_2(y)\}
\]
として表せば、積分順序を交換して
\[
\iint_{D} f(x,y)\,dA = \int_{c}^{d}\left(\int_{\psi_1(y)}^{\psi_2(y)} f(x,y)\,dx\right)dy
\]
としてもよい。両者は同じ重積分の値を与える。$x$ 積分が初等関数で書けない場合（例えば $e^{x^2}$）でも、順序を交換することで積分可能になることがある。

\bigskip
\noindent
\textbf{(2) 極座標変換とヤコビアン}

\noindent
$x = r\cos\theta,\ y = r\sin\theta$（$r \ge 0,\ 0 \le \theta < 2\pi$）と変数変換すると、ヤコビ行列は
\[
\frac{\partial(x,y)}{\partial(r,\theta)} =
\begin{pmatrix}
\dfrac{\partial x}{\partial r} & \dfrac{\partial x}{\partial \theta} \\[4pt]
\dfrac{\partial y}{\partial r} & \dfrac{\partial y}{\partial \theta}
\end{pmatrix}
=
\begin{pmatrix}
\cos\theta & -r\sin\theta \\
\sin\theta & r\cos\theta
\end{pmatrix}
\]
であり、そのヤコビアン（行列式）は
\[
J = \cos\theta \cdot r\cos\theta - (-r\sin\theta)\cdot\sin\theta = r(\cos^{2}\theta + \sin^{2}\theta) = r
\]
となる。したがって面積要素は $dx\,dy = |J|\,dr\,d\theta = r\,dr\,d\theta$ と変換され、
\[
\iint_{D} f(x,y)\,dx\,dy = \iint_{D'} f(r\cos\theta, r\sin\theta)\, r\,dr\,d\theta
\]
が成り立つ（$D'$ は $D$ を $(r,\theta)$ 平面に写した領域）。円板・扇形など、原点を中心とする領域を扱う際に有効である。

\bigskip
\noindent
\textbf{(3) 曲面下の体積}

\noindent
連続関数 $z = f(x,y) \ge 0$ のグラフと平面 $z=0$、および底面 $D$ とで囲まれる立体の体積は
\[
V = \iint_{D} f(x,y)\,dx\,dy
\]
で与えられる。これは、$D$ を微小面積 $dA$ に分割し、各微小面積の上に高さ $f(x,y)$ の柱を立てて体積 $f(x,y)\,dA$ を足し合わせる（リーマン和の極限をとる）ことから従う。回転対称な曲面（例えば $z = f(x^2+y^2)$）の場合は、極座標変換により積分が簡単になることが多い。

\clearpage
\begin{mondaiboxS}{問1}
次の累次積分の積分順序を交換し、その値を求めよ。
\end{mondaiboxS}
\[
\int_{0}^{1}\int_{y}^{1} e^{x^{2}}\,dx\,dy
\]
\kaitou
積分領域は $0 \le y \le 1,\ y \le x \le 1$ である。これは
\[
0 \le x \le 1,\qquad 0 \le y \le x
\]
と同じ領域を表す（$x$-$y$ 平面上で直線 $y=x$ の下側、$0\le x\le1$ の三角形領域）。したがって積分順序を交換すると
\[
\int_{0}^{1}\int_{y}^{1} e^{x^{2}}\,dx\,dy = \int_{0}^{1}\int_{0}^{x} e^{x^{2}}\,dy\,dx
\]
内側の積分は $x$ を定数とみて
\[
\int_{0}^{x} e^{x^{2}}\,dy = x\,e^{x^{2}}
\]
よって
\[
\int_{0}^{1} x\,e^{x^{2}}\,dx = \left[\frac{1}{2}e^{x^{2}}\right]_{0}^{1} = \frac{e-1}{2}
\]
となるから
\[
\boxed{\; \int_{0}^{1}\int_{y}^{1} e^{x^{2}}\,dx\,dy = \frac{e-1}{2} \;}
\]
（検算: $\dfrac{d}{dx}\left(\dfrac{1}{2}e^{x^2}\right) = x e^{x^2}$ なので原始関数の計算は正しい。また元の積分順序のままでは $\int e^{x^2}dx$ が初等関数で表せないため、順序交換が本質的である。）

\clearpage
\begin{mondaiboxS}{問2}
領域 $D = \{(x,y) \mid x^{2} + y^{2} \le 4,\ x \ge 0,\ y \ge 0\}$ において、次の重積分を極座標変換により求めよ。
\end{mondaiboxS}
\[
\iint_{D} \sqrt{x^{2} + y^{2}}\,dx\,dy
\]
\kaitou
$x = r\cos\theta,\ y = r\sin\theta$ とおくと、$D$ は第1象限の半径2の扇形なので
\[
D' = \{(r,\theta) \mid 0 \le r \le 2,\ 0 \le \theta \le \tfrac{\pi}{2}\}
\]
に対応する。$\sqrt{x^2+y^2} = r$、面積要素は $dx\,dy = r\,dr\,d\theta$ であるから
\[
\iint_{D} \sqrt{x^{2}+y^{2}}\,dx\,dy = \int_{0}^{\pi/2}\int_{0}^{2} r \cdot r\,dr\,d\theta
= \int_{0}^{\pi/2}\left(\int_{0}^{2} r^{2}\,dr\right)d\theta
\]
内側の積分は
\[
\int_{0}^{2} r^{2}\,dr = \left[\frac{r^{3}}{3}\right]_{0}^{2} = \frac{8}{3}
\]
であるから
\[
\int_{0}^{\pi/2} \frac{8}{3}\,d\theta = \frac{8}{3}\cdot\frac{\pi}{2} = \frac{4\pi}{3}
\]
よって
\[
\boxed{\; \iint_{D} \sqrt{x^{2}+y^{2}}\,dx\,dy = \frac{4\pi}{3} \;}
\]
（検算: 半径2、中心角 $\pi/2$ の扇形の「重み付き面積」として、$r=2$ での被積分関数の最大値は $2$、扇形の面積は $\pi\cdot2^2/4=\pi$ であるから、答え $4\pi/3 \approx 4.19$ は $0$ から $2\pi\approx6.28$ の間に収まり妥当である。）

\clearpage
\begin{mondaiboxS}{問3}
曲面 $z = 4 - x^{2} - y^{2}$ と平面 $z = 0$ とで囲まれる立体のうち、$z \ge 0$ の部分の体積 $V$ を求めよ。
\end{mondaiboxS}
\kaitou
$z \ge 0$ となるのは $x^{2}+y^{2} \le 4$ のときであるから、底面領域は
\[
D = \{(x,y) \mid x^{2}+y^{2} \le 4\}
\]
であり、求める体積は
\[
V = \iint_{D} (4 - x^{2} - y^{2})\,dx\,dy
\]
極座標 $x=r\cos\theta,\ y=r\sin\theta$（$0\le r\le 2,\ 0\le\theta\le 2\pi$）を用いると、$x^2+y^2=r^2$ であるから
\[
V = \int_{0}^{2\pi}\int_{0}^{2} (4 - r^{2})\, r\,dr\,d\theta
= \int_{0}^{2\pi}\int_{0}^{2} (4r - r^{3})\,dr\,d\theta
\]
内側の積分は
\[
\int_{0}^{2} (4r - r^{3})\,dr = \left[2r^{2} - \frac{r^{4}}{4}\right]_{0}^{2} = (8 - 4) - 0 = 4
\]
であるから
\[
V = \int_{0}^{2\pi} 4\,d\theta = 4 \cdot 2\pi = 8\pi
\]
よって
\[
\boxed{\; V = 8\pi \;}
\]
（検算: この立体は放物面 $z=4-r^2$（$0\le r\le2$）を底面とする回転体であり、円柱殻の方法でも
$V = \int_0^2 2\pi r (4-r^2)\,dr = 2\pi\left[2r^2 - \dfrac{r^4}{4}\right]_0^2 = 2\pi\cdot4 = 8\pi$
と一致し、答えが正しいことが確認できる。）

\end{document}
